\section{Инструкция по установке и запуску}

\subsection{Требования}

\begin{itemize}
    \item Java 17+ (JDK)
    \item Maven 3.8+
    \item Docker (для PostgreSQL и Testcontainers)
\end{itemize}

\subsection{Запуск приложения}

\begin{enumerate}
    \item Запустить PostgreSQL через Docker Compose:
\begin{verbatim}
docker-compose up -d
\end{verbatim}

    \item Собрать и запустить приложение:
\begin{verbatim}
mvn spring-boot:run
\end{verbatim}

    \item API доступно по адресу \texttt{http://localhost:8080}:
    \begin{itemize}
        \item \texttt{/api/orders} --- заказы (@RestController)
        \item \texttt{/api/payments} --- платежи (Functional Endpoints)
        \item \texttt{/api/products} --- товары (Spring Data REST)
        \item \texttt{/api/explorer} --- HAL Explorer (Spring Data REST)
    \end{itemize}
\end{enumerate}

\subsection{Запуск тестов}

\begin{enumerate}
    \item Запуск всех тестов (профиль с идемпотентностью):
\begin{verbatim}
mvn test
\end{verbatim}

    \item Запуск тестов без идемпотентности (демонстрация багов):
\begin{verbatim}
mvn test -Dtest="NonIdempotentProfileTest"
\end{verbatim}
\end{enumerate}

\textbf{Примечание:} для запуска тестов требуется Docker, так как Testcontainers автоматически запускает PostgreSQL в контейнере.

\subsection{Пример использования API}

Создание заказа с Idempotency-Key:
\begin{verbatim}
curl -X POST http://localhost:8080/api/orders \
  -H "Content-Type: application/json" \
  -H "Idempotency-Key: unique-key-123" \
  -d '{"description": "Ноутбук", "amount": 89999.99}'
\end{verbatim}

Повторный запрос с тем же ключом вернёт кэшированный ответ (заказ не дублируется).
